\documentclass[11pt,a4paper]{article}

\usepackage[utf8]{inputenc}
\usepackage[T1]{fontenc}
\usepackage[spanish]{babel}
\usepackage{amsmath,amssymb}
\usepackage[margin=2.5cm]{geometry}
\usepackage{hyperref}

\newcommand{\code}[1]{\texttt{#1}}
\emergencystretch=2em   % Evita que los nombres en \texttt se salgan del margen

\title{SpatFD v2: salvedades de la implementación}
\author{Brief técnico del \emph{pipeline} del paquete}
\date{Octubre de 2026}

\begin{document}
\maketitle

\noindent
Este documento resume las principales salvedades de la implementación actual
de SpatFD v2. No discute los métodos estadísticos en sí, sino las decisiones
del \emph{pipeline} del paquete: cómo se suaviza, cómo se hace el análisis de
componentes principales funcionales, cómo se traen los modelos de variograma
y, sobre todo, cómo se predice. Cada punto indica qué hace hoy el código y por
qué conviene tenerlo presente.

\section{Flujo del paquete}

El flujo completo tiene cuatro pasos:
\begin{enumerate}
  \item \code{spatfd()}: suavizamiento de las curvas y análisis de componentes
        principales funcionales (FPCA).
  \item \code{togstat()}: exportación de los \emph{scores} para estimar y
        ajustar variogramas por fuera del paquete.
  \item \code{add\_model()}: traducción, compilación y registro de los modelos
        ajustados dentro del objeto.
  \item \code{predict()}: predicción de las curvas en sitios nuevos o sobre un
        mapa, con los métodos \code{scores} y \code{lambda}.
\end{enumerate}
La predicción ya no depende de \code{gstat}: los sistemas se resuelven de
forma matricial dentro del paquete. Tanto el método \code{scores} como el
método \code{lambda} reproducen los resultados de la versión CRAN (0.0.1) con
diferencias relativas del orden de $10^{-16}$.

\section{Suavizamiento}

\begin{itemize}
  \item \textbf{Operador de penalización.} Se define con
        \code{vec2Lfd(c(1, ncol(dat)))}, heredado de la versión CRAN. Con $n$
        estaciones el operador es
        \[
          L\chi = D^2\chi + n\,D\chi + \chi ,
        \]
        de modo que la penalización depende del número de estaciones. No hay
        una justificación evidente para ello y cambia el suavizamiento cuando
        se agregan o quitan sitios.
  \item \textbf{Parámetros fijos.} El parámetro de suavizamiento
        \code{lambda} y el número de bases \code{nbasis} los elige el usuario
        y son los mismos para todas las curvas. No hay selección automática,
        por ejemplo por validación cruzada generalizada.
  \item \textbf{Datos faltantes.} No se revisan valores faltantes en los
        datos; la versión CRAN sí lo hacía.
\end{itemize}

\section{Componentes principales funcionales}

\begin{itemize}
  \item \textbf{Media constante en el espacio.} La función media $\mu(t)$ es
        el promedio simple de las curvas de las estaciones, sin considerar su
        ubicación. Todo lo que sigue supone que la media no cambia en el
        espacio.
  \item \textbf{Dependencia del número de componentes.} Los \emph{scores}
        cambian ligeramente, en el cuarto decimal, según cuántos componentes
        se pidan, por la integración numérica de \code{fda}. Por eso, cuando
        se selecciona el número de componentes por varianza explicada, el
        análisis se repite en lugar de recortarse.
  \item \textbf{Un FPCA por variable.} Cada variable funcional tiene su
        propio análisis; no hay un análisis conjunto entre variables.
\end{itemize}

\section{Modelamiento de los variogramas}

\begin{itemize}
  \item \textbf{Pocas estaciones.} Con pocas estaciones (13 para PM10 en los
        datos de México) los variogramas empíricos son muy inestables. En el
        ejemplo de demostración el ajuste de los variogramas de PM10 no
        converge.
  \item \textbf{Meseta frente a valor propio.} Nada obliga a que la meseta
        del modelo ajustado se parezca al valor propio del componente, que es
        la varianza empírica del \emph{score}. La diferencia entre ambos
        afecta directamente la varianza de predicción.
  \item \textbf{Traducción desde gstat.} \code{toSpatFD()} solo conserva el
        modelo, el rango, la meseta parcial y el parámetro \code{kappa}. La
        anisotropía se descarta sin aviso, y un modelo \code{Lin} con rango se
        traduce como lineal no acotado.
  \item \textbf{Modelos espejo de gstat.} Los 18 modelos espejo siguen las
        fórmulas del código fuente de \code{gstat} y reproducen sus valores.
        Dos heredan particularidades de \code{gstat}: \code{spline} devuelve
        la meseta más allá de su rango (en un modelo anidado \code{gstat}
        aplica ese corte en el mayor de los rangos, aquí en el rango de la
        propia estructura), y \code{legendre} toma las distancias en
        kilómetros sobre la esfera.
  \item \textbf{Un solo traductor.} Solo se pueden traer modelos ajustados con
        \code{gstat}.
\end{itemize}

\section{Predicción}

\subsection{Aspectos comunes a los dos métodos}

\begin{itemize}
  \item \textbf{Kriging simple con media cero.} Las covarianzas se obtienen de
        los variogramas como $C(h) = C(0) - \gamma(h)$, con $C(0)$ la suma de
        las mesetas parciales. Esto descansa en que los \emph{scores} sean
        estacionarios y centrados.
  \item \textbf{Modelos sin meseta.} Los modelos \code{linear}, \code{power}
        y \code{logarithmic} no tienen covarianza y no se pueden usar hasta
        que exista kriging ordinario.
  \item \textbf{Distancias euclidianas.} Las distancias se calculan sobre las
        coordenadas tal como vienen. Sirven para coordenadas proyectadas, no
        para longitud y latitud.
  \item \textbf{Vecindario global.} Se factoriza una matriz densa con todos
        los sitios observados, lo que escala mal cuando hay muchos sitios. No
        hay kriging por vecindarios.
  \item \textbf{Validez del modelo.} Solo se detecta un modelo inválido si la
        factorización de Cholesky falla. Un modelo que no sea válido puede
        pasar sin error.
  \item \textbf{Efecto pepita.} Se trata como en \code{gstat}: el variograma
        vale cero en el origen y la predicción en un sitio observado devuelve
        el valor observado.
\end{itemize}

\subsection{Método \code{scores}}

Cada \emph{score} se predice por kriging simple y la curva se reconstruye:
\[
  \hat f_k(s_0) = c_0^{T}\,\Sigma^{-1} f,
  \qquad
  \hat\chi_{s_0}(t) = \mu(t) + \sum_{k=1}^{K} \hat f_k(s_0)\,\xi_k(t).
\]
\begin{itemize}
  \item \textbf{Independencia entre scores.} Los \emph{scores} de una misma
        variable se tratan como independientes salvo que el modelo traiga
        variogramas cruzados. El FPCA solo garantiza que no estén
        correlacionados en un mismo sitio, no a distancia.
  \item \textbf{Truncamiento.} La curva se reconstruye con los $K$ componentes
        retenidos, y la varianza reportada, $\sum_k \sigma_k^2(s_0)$, ignora la
        variabilidad de los componentes que quedaron por fuera.
  \item \textbf{Un solo sistema.} Los casos univariado y multivariado pasan
        por el mismo \emph{solver} por bloques. Los \emph{scores} que no están
        conectados por variogramas cruzados se resuelven por aparte.
\end{itemize}

\subsection{Método \code{lambda}}

Se implementó igual que en la versión CRAN, por decisión explícita, para
permitir la comparación directa:
\[
  \hat\chi_{s_0}(t) = \sum_{i=1}^{n} \lambda_i\,\chi_{s_i}(t),
  \qquad
  \Omega\,\lambda = c,
  \qquad
  \Omega = \sum_{k=1}^{K}\Omega_k,\quad c = \sum_{k=1}^{K} c_k ,
\]
con varianza de predicción
\[
  \sum_{k=1}^{K}\left(\eta_k - 2\,c_k^{T}\lambda + \lambda^{T}\Omega_k\lambda\right),
\]
donde $\eta_k$ es el valor propio del componente $k$.
\begin{itemize}
  \item \textbf{Curvas sin centrar.} Los pesos son de kriging simple y no
        suman uno, pero se aplican a las curvas observadas sin restarles la
        media. La predicción se encoge hacia cero y no hacia la media: en PM10
        las curvas predichas quedan muy por debajo de las observadas. La
        versión coherente sería
        \[
          \hat\chi_{s_0}(t) = \mu(t) + \sum_{i=1}^{n}\lambda_i
          \bigl(\chi_{s_i}(t) - \mu(t)\bigr),
        \]
        y el lugar donde hacer el cambio está marcado en el código de
        \code{predict\_lambda()}.
  \item \textbf{Varianza con el valor propio.} La varianza mezcla el valor
        propio $\eta_k$ con las covarianzas del modelo ajustado. Si la meseta
        y el valor propio difieren mucho puede salir negativa, y no hay aviso.
  \item \textbf{No comparable con \code{scores}.} Por las dos razones
        anteriores, ni las curvas ni las varianzas del método \code{lambda}
        son comparables hoy con las del método \code{scores}.
  \item \textbf{Solo univariado.} Si el modelo tiene variogramas cruzados o se
        pide más de una variable, la función se detiene con un error.
\end{itemize}

\section{Objeto y registro de modelos}

\begin{itemize}
  \item \textbf{Registro en el sitio.} \code{add\_model()} modifica el objeto
        sin reasignarlo. Una copia hecha con \code{\textless-} comparte los modelos
        del objeto original.
  \item \textbf{Identificadores de variogramas cruzados.} Se arman uniendo los
        nombres de los \emph{scores} con un punto. Nombres de variable que
        contengan puntos pueden volverse ambiguos.
  \item \textbf{Mapas.} Al predecir sobre un mapa no se revisa ni se
        transforma el sistema de coordenadas: debe ser el mismo de las
        estaciones. El mapa de predicción muestra un solo instante.
\end{itemize}

\section{Validación cruzada}

\code{leave\_one\_out()} compara modelos y métodos dejando por fuera un sitio
a la vez y midiendo la distancia entre la curva observada y la predicha,
$\lVert \chi_{s_i} - \hat\chi_{s_i}^{(-i)} \rVert$.
\begin{itemize}
  \item \textbf{Qué se repite y qué no.} Sin el sitio se repiten el
        suavizamiento y el FPCA, con los mismos argumentos y el mismo número
        de componentes. Los modelos de variograma \emph{no} se vuelven a
        ajustar: son los registrados en el objeto, que se ajustaron con todos
        los sitios. El error queda por eso algo optimista.
  \item \textbf{Signo de los componentes.} El signo de un componente
        principal es arbitrario y puede cambiar al quitar un sitio. Los
        componentes de cada repetición se dejan con el signo del análisis
        completo, para que los variogramas cruzados conserven su sentido.
  \item \textbf{Cokriging.} Con un modelo de varias variables solo se deja por
        fuera la curva de la variable validada; las demás variables conservan
        su observación en ese sitio.
  \item \textbf{Diferencia con CRAN.} Las curvas predichas en cada repetición
        son las mismas que las de \code{crossval\_loo()} de la versión CRAN.
        Los errores difieren hasta en un 2\% porque CRAN calcula la norma
        por integración numérica y aquí se calcula de forma exacta con los
        productos internos de la base.
  \item \textbf{Costo.} Cada sitio implica repetir el análisis y factorizar de
        nuevo las matrices de covarianza, por cada modelo comparado.
\end{itemize}

\section{Pendientes}

\begin{itemize}
  \item Kriging ordinario, necesario para los modelos sin meseta.
  \item Método \code{lambda} centrado y su extensión multivariada.
\end{itemize}

\end{document}
